\section{2015 Probability and Statistics}

\subsection{Individual}
Please solve the following 5 problems.

\begin{exercise}
(a) Let $X$ and $Y$ be two random variables with zero means, variance 1, and correlation $\rho$. Prove that
\[
\mathbb{E}[\max\{X^2, Y^2\}] \leq 1 + \sqrt{1-\rho^2}.
\]

(b) Let $X$ and $Y$ have a bivariate normal distribution with zero means, variances $\sigma^2$ and $\tau^2$, respectively, and correlation $\rho$. Find the conditional expectation $\mathbb{E}(X | Y)$.
\end{exercise}

\begin{solution}
(a) We use the identity $\max\{a, b\} = \frac{a+b+|a-b|}{2}$. Thus,
\[
\mathbb{E}[\max\{X^2, Y^2\}] = \frac{\mathbb{E}[X^2] + \mathbb{E}[Y^2] + \mathbb{E}[|X^2 - Y^2|]}{2} = 1 + \frac{1}{2}\mathbb{E}[|X^2 - Y^2|].
\]
Since $X^2 - Y^2 = (X-Y)(X+Y)$, by the Cauchy-Schwarz inequality, we have
\[
\mathbb{E}[|X^2 - Y^2|] \leq \sqrt{\mathbb{E}[(X-Y)^2] \mathbb{E}[(X+Y)^2]}.
\]
Given that $\mathbb{E}[X] = \mathbb{E}[Y] = 0$, $\text{Var}(X) = \text{Var}(Y) = 1$, and $\text{Corr}(X, Y) = \rho$, we find
\[
\mathbb{E}[(X-Y)^2] = \text{Var}(X-Y) = \text{Var}(X) + \text{Var}(Y) - 2\text{Cov}(X, Y) = 1 + 1 - 2\rho = 2(1-\rho),
\]
\[
\mathbb{E}[(X+Y)^2] = \text{Var}(X+Y) = \text{Var}(X) + \text{Var}(Y) + 2\text{Cov}(X, Y) = 1 + 1 + 2\rho = 2(1+\rho).
\]
Therefore,
\[
\mathbb{E}[|X^2 - Y^2|] \leq \sqrt{2(1-\rho) \cdot 2(1+\rho)} = 2\sqrt{1-\rho^2}.
\]
Plugging this back into the expression for the expectation:
\[
\mathbb{E}[\max\{X^2, Y^2\}] \leq 1 + \frac{1}{2} \cdot 2\sqrt{1-\rho^2} = 1 + \sqrt{1-\rho^2}.
\]

(b) For a bivariate normal distribution $(X, Y) \sim N\left(\begin{pmatrix} 0 \\ 0 \end{pmatrix}, \begin{pmatrix} \sigma^2 & \rho\sigma\tau \\ \rho\sigma\tau & \tau^2 \end{pmatrix}\right)$, the conditional distribution of $X$ given $Y=y$ is normal with mean
\[
\mathbb{E}[X|Y=y] = \mu_X + \frac{\text{Cov}(X, Y)}{\text{Var}(Y)}(y - \mu_Y).
\]
Substituting the given parameters:
\[
\mathbb{E}[X|Y] = 0 + \frac{\rho\sigma\tau}{\tau^2}(Y - 0) = \rho \frac{\sigma}{\tau} Y.
\]
\end{solution}

\begin{exercise}
We flip a fair coin until heads turns out twice consecutively. What is the expected number of flips?
\end{exercise}

\begin{solution}
Let $E$ be the expected number of additional flips needed to see two consecutive heads, starting from the beginning.
Let $E_0$ be the expected number of additional flips when we need two more heads.
Let $E_1$ be the expected number of additional flips when we have just flipped one head and need one more.
In the initial state (state 0), we flip the coin:
- With probability $1/2$, we get tails (T) and stay in state 0.
- With probability $1/2$, we get heads (H) and move to state 1.
Thus, $E_0 = 1 + \frac{1}{2}E_0 + \frac{1}{2}E_1 \implies \frac{1}{2}E_0 = 1 + \frac{1}{2}E_1 \implies E_0 = 2 + E_1$.
In state 1:
- With probability $1/2$, we get tails (T) and go back to state 0.
- With probability $1/2$, we get heads (H) and are done (0 more flips needed).
Thus, $E_1 = 1 + \frac{1}{2}E_0 + \frac{1}{2}(0) = 1 + \frac{1}{2}E_0$.
Substituting $E_1$ into the equation for $E_0$:
\[
E_0 = 2 + (1 + \frac{1}{2}E_0) = 3 + \frac{1}{2}E_0.
\]
Solving for $E_0$:
\[
\frac{1}{2}E_0 = 3 \implies E_0 = 6.
\]
The expected number of flips is 6.
\end{solution}

\begin{exercise}
Let $(X_n, n \geq 1)$ be a sequence of independent Gaussian variables, with respective mean $\mu_n$, and variance $\sigma_n^2$.

(a) Prove that if $\sum_n X_n^2$ converges in $L^1$, then $\sum_n X_n^2$ converges in $L^p$, for every $p \in [1, \infty)$.

(b) Assume that $\mu_n = 0$, for every $n$. Prove that if $\sum_n \sigma_n^2 = \infty$, then
\[
\mathbb{P}\left(\sum_n X_n^2 = \infty\right) = 1.
\]
\end{exercise}

\begin{solution}
(a) Let $S_N = \sum_{n=1}^N X_n^2$. Since $X_n^2 \geq 0$, $S_N$ is non-decreasing. $L^1$ convergence implies that $\mathbb{E}[\sum_n X_n^2] = \sum_n \mathbb{E}[X_n^2] = \sum_n (\mu_n^2 + \sigma_n^2) < \infty$.
Let $S = \sum_n X_n^2$. For independent Gaussian $X_n$, $X_n^2$ has a non-central chi-squared distribution. The moment generating function of $X_n^2$ is $\mathbb{E}[e^{t X_n^2}] = \frac{1}{\sqrt{1-2t\sigma_n^2}} \exp\left(\frac{t\mu_n^2}{1-2t\sigma_n^2}\right)$ for $t < \frac{1}{2\sigma_n^2}$.
Since $\sigma_n^2 \to 0$, for any $t > 0$, there exists $N_t$ such that $2t\sigma_n^2 < 1$ for all $n > N_t$. The MGF of $S$ is $\prod_n \mathbb{E}[e^{t X_n^2}]$. For small $t$, the logarithm of the MGF is $\sum_n \left[ -\frac{1}{2} \ln(1-2t\sigma_n^2) + \frac{t\mu_n^2}{1-2t\sigma_n^2} \right]$.
Using $\ln(1-x) \approx -x$, the terms are approximately $t\sigma_n^2 + t\mu_n^2 = t(\sigma_n^2 + \mu_n^2)$. Since $\sum (\sigma_n^2 + \mu_n^2) < \infty$, the sum of logarithms converges for small $t > 0$. Thus, $S$ has a finite MGF in a neighborhood of 0, which implies all moments $\mathbb{E}[S^p]$ are finite for $p \in [1, \infty)$. Monotone convergence then ensures $L^p$ convergence.

(b) If $\mu_n = 0$, then $X_n \sim N(0, \sigma_n^2)$, so $X_n^2 = \sigma_n^2 Z_n^2$ where $Z_n^2 \sim \chi_1^2$.
The sum $S = \sum_n X_n^2$ consists of independent non-negative random variables. By Kolmogorov's Three-Series Theorem (or a simpler version for non-negative variables), $\sum Y_n$ diverges almost surely if $\sum \mathbb{E}[Y_n \wedge c] = \infty$ for some $c > 0$.
Here $\mathbb{E}[X_n^2 \wedge 1] = \sigma_n^2 \mathbb{E}[Z_n^2 \wedge \frac{1}{\sigma_n^2}]$. As $n \to \infty$, if $\sigma_n^2 \to 0$, $\mathbb{E}[Z_n^2 \wedge \frac{1}{\sigma_n^2}] \to \mathbb{E}[Z_n^2] = 1$. Thus $\mathbb{E}[X_n^2 \wedge 1] \sim \sigma_n^2$. Since $\sum \sigma_n^2 = \infty$, the sum of expectations of truncated variables diverges, so $\sum X_n^2 = \infty$ almost surely. If $\sigma_n^2 \not\to 0$, divergence is even more trivial.
\end{solution}

\begin{exercise}
Let $X_1, \ldots, X_n$ be a random sample of size $n$ from the exponential distribution with pdf $f(x; \theta) = \theta^{-1}\exp(-x/\theta)$ for $x, \theta > 0$, zero elsewhere. Let $Y_1 = \min\{X_1, \dots, X_n\}$. Consider an estimator $nY_1$.

(a) Show this estimate is unbiased.

(b) Prove or disprove: This estimate is a consistent estimator.

(c) Prove or disprove: This estimate is an efficient estimator.
\end{exercise}

\begin{solution}
(a) The distribution of the minimum $Y_1$ of $n$ i.i.d. $\text{Exp}(\theta)$ variables is $\text{Exp}(\theta/n)$. The expectation is $\mathbb{E}[Y_1] = \theta/n$.
Therefore, $\mathbb{E}[nY_1] = n \mathbb{E}[Y_1] = n(\theta/n) = \theta$. The estimator is unbiased.

(b) Disprove. For a consistent estimator $\hat{\theta}_n$, we require $\hat{\theta}_n \xrightarrow{P} \theta$ as $n \to \infty$.
However, the distribution of $nY_1$ is $\text{Exp}(\theta)$ for all $n$. It does not depend on $n$ and does not degenerate to the constant $\theta$. Thus, $nY_1$ does not converge in probability to $\theta$.

(c) Disprove. An efficient estimator achieves the Cram\'er-Rao Lower Bound (CRLB).
The Fisher information for a sample of size $n$ from $\text{Exp}(\theta)$ is $I_n(\theta) = n/\theta^2$. The CRLB is $1/I_n(\theta) = \theta^2/n$.
The variance of our estimator is $\text{Var}(nY_1) = n^2 \text{Var}(Y_1) = n^2 (\theta/n)^2 = \theta^2$.
For $n > 1$, $\theta^2 > \theta^2/n$, so the estimator is not efficient.
\end{solution}

\begin{exercise}
Let the independent normal random variables $Y_1, \dots, Y_n$ have, respectively, the probability density functions $N(\mu, \gamma^2 x_i^2)$, $i = 1, \ldots, n$, where the given $x_1, \ldots, x_n$ are not all equal and no one of which is zero.

(a) Construct a confidence interval for $\gamma$ with significance level $1-\alpha$.

(b) Discuss the test of the hypothesis $H_0: \gamma = 1$, $\mu$ unspecified, against all alternatives $H_1: \gamma \neq 1$, $\mu$ unspecified.
\end{exercise}

\begin{solution}
(a) Let $Z_i = Y_i/x_i$. Then $Z_i \sim N(\mu/x_i, \gamma^2)$. This is a linear model $Z_i = \mu w_i + \epsilon_i$ with $w_i = 1/x_i$ and $\epsilon_i \sim N(0, \gamma^2)$.
The least squares estimator for $\mu$ is $\hat{\mu} = \frac{\sum Z_i w_i}{\sum w_i^2} = \frac{\sum Y_i/x_i^2}{\sum 1/x_i^2}$.
The residual sum of squares is $RSS = \sum (Z_i - \hat{\mu} w_i)^2 = \sum \frac{(Y_i - \hat{\mu})^2}{x_i^2}$.
It is a standard result that $Q = \frac{RSS}{\gamma^2} \sim \chi_{n-1}^2$.
A $1-\alpha$ confidence interval for $\gamma^2$ is $\left[ \frac{RSS}{\chi_{n-1, \alpha/2}^2}, \frac{RSS}{\chi_{n-1, 1-\alpha/2}^2} \right]$.
Taking the square root gives the interval for $\gamma$:
\[
\left[ \sqrt{\frac{\sum (Y_i - \hat{\mu})^2/x_i^2}{\chi_{n-1, \alpha/2}^2}}, \sqrt{\frac{\sum (Y_i - \hat{\mu})^2/x_i^2}{\chi_{n-1, 1-\alpha/2}^2}} \right].
\]

(b) To test $H_0: \gamma=1$ against $H_1: \gamma \neq 1$, we use the test statistic $Q_0 = \sum \frac{(Y_i - \hat{\mu})^2}{x_i^2}$.
Under $H_0$, $Q_0 \sim \chi_{n-1}^2$.
At significance level $\alpha$, we reject $H_0$ if $Q_0 > \chi_{n-1, \alpha/2}^2$ or $Q_0 < \chi_{n-1, 1-\alpha/2}^2$.
This is a two-tailed test based on the chi-squared distribution.
\end{solution}

\subsection{Team}
Please solve the following 5 problems.

\begin{exercise}
One hundred passengers board a plane with exactly 100 seats. The first passenger takes a seat at random. The second passenger takes his own seat if it is available, otherwise he takes at random a seat among the available ones. The third passenger takes his own seat if it is available, otherwise he takes at random a seat among the available ones. This process continues until all the 100 passengers have boarded the plane. What is the probability that the last passenger takes his own seat?
\end{exercise}

\begin{solution}
Consider the moment when a passenger is forced to choose a seat at random (either the first passenger or a subsequent one whose seat was taken). They will either choose seat 1, seat 100, or some other seat $k$.
If they choose seat 1, then passenger 100 will eventually find seat 100 free.
If they choose seat 100, then passenger 100 will definitely not get his own seat.
If they choose some seat $k$ ($1 < k < 100$), then passengers $2, \dots, k-1$ will all find their own seats. Passenger $k$ will then be the one to choose a seat at random.
Crucially, at any such choice, seat 1 and seat 100 are equally likely to be picked. The first of these two seats to be picked determines the fate of passenger 100. If seat 1 is picked first, passenger 100 gets his seat. If seat 100 is picked first, he doesn't.
By symmetry, the probability is $1/2$.
\end{solution}

\begin{exercise}
Assume a sequence of random variables $X_n$ converges in distribution to a random variable $X$. Let $\{N_t, t \geq 0\}$ be a set of positive integer-valued random variables, which is independent of $(X_n)$ and converges in probability to $\infty$ as $t \to \infty$. Prove that $X_{N_t}$ converges in distribution to $X$ as $t \to \infty$.
\end{exercise}

\begin{solution}
Let $F_n$ be the CDF of $X_n$ and $F$ be the CDF of $X$. We want to show $P(X_{N_t} \leq x) \to F(x)$ at all continuity points $x$ of $F$.
By independence, $P(X_{N_t} \leq x) = \sum_{n=1}^\infty P(X_n \leq x) P(N_t = n) = \mathbb{E}[F_{N_t}(x)]$.
Since $X_n \xrightarrow{d} X$, $F_n(x) \to F(x)$ for continuity points $x$.
For any $\epsilon > 0$, there exists $M$ such that $|F_n(x) - F(x)| < \epsilon$ for all $n > M$.
Then
\[
|\mathbb{E}[F_{N_t}(x)] - F(x)| \leq \mathbb{E}[|F_{N_t}(x) - F(x)|]
\]
\[
= \sum_{n \leq M} |F_n(x) - F(x)| P(N_t = n) + \sum_{n > M} |F_n(x) - F(x)| P(N_t = n)
\]
\[
\leq 1 \cdot P(N_t \leq M) + \epsilon \cdot P(N_t > M).
\]
As $t \to \infty$, $P(N_t \leq M) \to 0$ because $N_t \xrightarrow{P} \infty$.
Thus $\limsup_{t \to \infty} |\mathbb{E}[F_{N_t}(x)] - F(x)| \leq \epsilon$. Since $\epsilon$ is arbitrary, the result follows.
\end{solution}

\begin{exercise}
Suppose $T_1, T_2, \dots, T_n$ is a sequence of i.i.d. random variables with the exponential distribution of the density function
\[
p(x) = \begin{cases} e^{-x}, & x \geq 0; \\ 0, & x < 0. \end{cases}
\]
Let $S_n = T_1 + T_2 + \cdots + T_n$. Find the distribution of the random vector
\[
V_n = \left(\frac{T_1}{S_n}, \frac{T_2}{S_n}, \dots, \frac{T_n}{S_n}\right).
\]
\end{exercise}

\begin{solution}
The joint density of $(T_1, \dots, T_n)$ is $f(t_1, \dots, t_n) = e^{-\sum t_i}$ for $t_i \geq 0$.
Consider the transformation $Y_i = T_i/S_n$ for $i=1, \dots, n-1$ and $Y_n = S_n$.
The inverse transformation is $T_i = Y_i Y_n$ for $i < n$ and $T_n = Y_n(1 - \sum_{i=1}^{n-1} Y_i)$.
The Jacobian matrix $J$ has elements $\frac{\partial T_i}{\partial Y_j}$.
One can compute $|\det J| = Y_n^{n-1}$.
The joint density of $(Y_1, \dots, Y_n)$ is
\[
g(y_1, \dots, y_n) = e^{-y_n} y_n^{n-1}, \quad y_n > 0, y_i > 0, \sum_{i=1}^{n-1} y_i < 1.
\]
The marginal density of $(Y_1, \dots, Y_{n-1})$ is
\[
h(y_1, \dots, y_{n-1}) = \int_0^\infty y_n^{n-1} e^{-y_n} dy_n = \Gamma(n) = (n-1)!.
\]
This is a constant value on the simplex $\Delta_{n-1} = \{ (y_1, \dots, y_{n-1}) : y_i > 0, \sum y_i < 1 \}$.
Thus, $V_n$ follows the Dirichlet distribution $\text{Dir}(1, 1, \dots, 1)$, which is the uniform distribution on the $(n-1)$-dimensional unit simplex.
\end{solution}

\begin{exercise}
Suppose that $X$ and $Z$ are jointly normal with mean zero and standard deviation 1. For a strictly monotonic function $f(\cdot)$, $\mathrm{cov}(X, Z) = 0$ if and only if $\mathrm{cov}(X, f(Z)) = 0$, provided the latter covariance exists.

Hint: $Z$ can be expressed as $Z = \rho X + \varepsilon$ where $X$ and $\varepsilon$ are independent and $\varepsilon \sim N(0, \sqrt{1-\rho^2})$.
\end{exercise}

\begin{solution}
If $\text{Cov}(X, Z) = 0$, then for jointly normal variables, $X$ and $Z$ are independent. Thus $X$ and $f(Z)$ are independent, and $\text{Cov}(X, f(Z)) = 0$.
Conversely, assume $\text{Cov}(X, f(Z)) = 0$. Let $Z = \rho X + \sqrt{1-\rho^2} W$ where $W \sim N(0,1)$ is independent of $X$.
$\text{Cov}(X, f(Z)) = \mathbb{E}[X f(\rho X + \sqrt{1-\rho^2} W)] = 0$.
Let $g(u) = \mathbb{E}_W[f(u + \sqrt{1-\rho^2} W)]$. Then $\mathbb{E}[X g(\rho X)] = 0$.
If $f$ is strictly increasing, then $g$ is strictly increasing (the expectation of an increasing function of a shifted Gaussian is increasing in the shift).
Since $\mathbb{E}[X] = 0$, $\text{Cov}(X, g(\rho X)) = \mathbb{E}[X g(\rho X)] - \mathbb{E}[X]\mathbb{E}[g(\rho X)] = \mathbb{E}[X g(\rho X)] = 0$.
For any random variable $X$ and strictly increasing function $h$, $\text{Cov}(X, h(X)) \geq 0$, and the covariance is zero if and only if $X$ is constant or $h$ is constant.
Here $g(\rho X)$ is strictly increasing in $X$ if $\rho > 0$ and strictly decreasing if $\rho < 0$.
If $\rho > 0$, $\text{Cov}(X, g(\rho X)) > 0$. If $\rho < 0$, $\text{Cov}(X, g(\rho X)) < 0$.
The only way the covariance can be zero is if $\rho = 0$.
Thus $\text{Cov}(X, Z) = \rho = 0$. The case for strictly decreasing $f$ follows by considering $-f$.
\end{solution}

\begin{exercise}
Consider the following penalized least-squares problem (Lasso):
\[
\frac{1}{2}\|\mathbf{Y} - \mathbf{X}\boldsymbol{\beta}\|^2 + \lambda\|\boldsymbol{\beta}\|_1.
\]

Let $\widehat{\beta}$ be a minimizer and $\Delta = \widehat{\beta} - \beta^*$ for any given $\beta^*$. If $\lambda > 2\|\mathbf{X}^T(\mathbf{Y} - \mathbf{X}\beta^*)\|_\infty$ show that:

1. $\|\mathbf{Y} - \mathbf{X}^T\widehat{\beta}\|^2 - \|\mathbf{Y} - \mathbf{X}^T\beta^*\|^2 > -\lambda\|\Delta\|_1$.

2. $\|\Delta_{S^c}\|_1 \leq 3\|\Delta_S\|_1$, where $S = \{j: \beta_j^* \neq 0\}$ is the support of the vector $\beta^*$, $S^c$ is its complement set, $\Delta_S$ is the subvector of $\Delta$ restricted on the set $S$, and $\|\Delta_S\|_1$ is its $L_1$-norm.
\end{exercise}

\begin{solution}
1. Let $f(\beta) = \frac{1}{2}\|\mathbf{Y} - \mathbf{X}\beta\|^2$. By convexity, $f(\widehat{\beta}) \geq f(\beta^*) + \langle \nabla f(\beta^*), \widehat{\beta} - \beta^* \rangle$.
$\nabla f(\beta^*) = -\mathbf{X}^T(\mathbf{Y} - \mathbf{X}\beta^*)$.
Thus, $f(\widehat{\beta}) - f(\beta^*) \geq -\langle \mathbf{X}^T(\mathbf{Y} - \mathbf{X}\beta^*), \Delta \rangle$.
By H\"older's inequality, $|\langle \mathbf{X}^T(\mathbf{Y} - \mathbf{X}\beta^*), \Delta \rangle| \leq \|\mathbf{X}^T(\mathbf{Y} - \mathbf{X}\beta^*)\|_\infty \|\Delta\|_1$.
Given $\lambda > 2\|\mathbf{X}^T(\mathbf{Y} - \mathbf{X}\beta^*)\|_\infty$, we have $\|\mathbf{X}^T(\mathbf{Y} - \mathbf{X}\beta^*)\|_\infty < \lambda/2$.
Therefore, $f(\widehat{\beta}) - f(\beta^*) > -\frac{\lambda}{2} \|\Delta\|_1 > -\lambda \|\Delta\|_1$.
Multiplying by 2 gives the result for $\|\mathbf{Y} - \mathbf{X}\beta\|^2$.

2. From the definition of $\widehat{\beta}$:
$f(\widehat{\beta}) + \lambda\|\widehat{\beta}\|_1 \leq f(\beta^*) + \lambda\|\beta^*\|_1$.
$f(\widehat{\beta}) - f(\beta^*) \leq \lambda(\|\beta^*\|_1 - \|\widehat{\beta}\|_1)$.
Using the lower bound from part 1:
$-\frac{\lambda}{2}\|\Delta\|_1 < \lambda(\|\beta^*\|_1 - \|\widehat{\beta}\|_1) \implies - \frac{1}{2}(\|\Delta_S\|_1 + \|\Delta_{S^c}\|_1) < \|\beta^*_S\|_1 - \|\widehat{\beta}_S\|_1 - \|\widehat{\beta}_{S^c}\|_1$.
Since $\beta^*_{S^c} = 0$, $\|\widehat{\beta}_{S^c}\|_1 = \|\Delta_{S^c}\|_1$. Also $\|\beta^*_S\|_1 - \|\widehat{\beta}_S\|_1 \leq \|\beta^*_S - \widehat{\beta}_S\|_1 = \|\Delta_S\|_1$.
$- \frac{1}{2}\|\Delta_S\|_1 - \frac{1}{2}\|\Delta_{S^c}\|_1 < \|\Delta_S\|_1 - \|\Delta_{S^c}\|_1$.
Rearranging terms:
$\frac{1}{2}\|\Delta_{S^c}\|_1 < \frac{3}{2}\|\Delta_S\|_1 \implies \|\Delta_{S^c}\|_1 < 3\|\Delta_S\|_1$.
(Note: The non-strict inequality $\|\Delta_{S^c}\|_1 \leq 3\|\Delta_S\|_1$ is thus also proven).
\end{solution}
